\section{Discussion}
\label{sec:discussion}

We motivate ICD and AICD through two distinct mechanisms, and then note the
assumptions and costs that the construction carries.

\subsection{ICD as model-adaptive difficulty}
\label{sec:discussion:icd}

ICD removes exactly the evidence the model currently finds most convincing.
Consider a classifier that has settled on a single dominant cue, such as a dog's
face or a flower's centre.  Under ICD it repeatedly encounters that cue occluded
and must predict from what remains.  This is a form of difficulty that adapts to the
model rather than to the data: the same image presents a harder problem to a
network that relies narrowly on one region than to one that already integrates
several.  The mechanism is related to the motivation of
curriculum learning~\citep{bengio2009curriculum}, with the ordering inverted.
Where a curriculum presents easy examples first and hard examples later, ICD
holds the data fixed and instead removes, at each step, whatever the model has
found easy.  The target it removes therefore moves, defined by the model's own
saliency.  Because the
saliency in Eq.~\eqref{eq:gradcam} is recomputed from the current parameters,
the occlusion follows the representation as it develops rather than being fixed
in advance.

\subsection{AICD as background suppression}
\label{sec:discussion:aicd}

AICD addresses the complementary failure mode.  In many natural datasets the
background is statistically informative: water co-occurs with certain animals,
indoor scenes with certain objects~\citep{beery2018recognition}.  A model can
exploit this without ever attending to the object, and such reliance is
invisible to validation accuracy when the validation set shares the same
background statistics.  AICD masks precisely the low-saliency region, which is
the background once the model has localised the object, and perturbs it
with a fill.  Repeatedly altering the background while preserving the object
discourages the classifier from treating background as evidence, in the spirit
of work that removes or randomises background to test and improve object-level
reliance~\citep{xiao2021noise}.

\subsection{Expected effect on saliency}
\label{sec:discussion:saliency}

Both methods act on the saliency map and may, in turn, reshape it.  Suppressing
the high-saliency region (ICD) should discourage the concentration of saliency
on any single area, so we would expect a model trained with ICD to develop a
more distributed saliency map; suppressing the background (AICD) should, if it
succeeds, sharpen saliency onto the object.  These expectations are stated here
as hypotheses.  Quantitative measures such as the fraction of saliency mass on
the image border, the coverage above a threshold, and the Gini concentration of
the map can make them testable, and we take up that measurement in the companion
study rather than claiming any outcome here.

\subsection{Assumptions and limitations}
\label{sec:discussion:limits}

The construction rests on several assumptions that bound its applicability.

\paragraph{Saliency must be meaningful.}
ICD and AICD are only as good as the saliency map that drives them.  At
initialisation, before the network has learned useful features, the GradCAM map
is close to uninformative, and masking by it approaches random masking.  The
methods are therefore best applied once the model has partially trained, or with
saliency computed from a warmed-up checkpoint; applying them from the first step
is unlikely to differ from Cutout.

\paragraph{Online cost.}
Unlike dataset-level augmentations that can be precomputed, the saliency in
Eq.~\eqref{eq:gradcam} depends on the current parameters and in principle must
be recomputed as training proceeds.  A gradient evaluation per image per update
is expensive.  Section~\ref{sec:implementation} describes the caching strategy
we use to amortise this, but the cache trades freshness for speed: a cached map
reflects the model at the time of caching, not the current step.

\paragraph{Tile granularity.}
The tile size $t$ is a fixed grid imposed on a continuous saliency map.  A small
$t$ can mask fine object parts that happen to be salient; a large $t$ can mask
substantial background along with the object.  The appropriate $t$ depends on
the object scale relative to the image, which the method does not adapt
automatically.

\paragraph{Single-image, single-label scope.}
ICD and AICD operate on one image and leave the label unchanged, which keeps
them simple and composable but also means they do not exploit the label-mixing
that benefits methods such as CutMix~\citep{yun2019cutmix} and
SnapMix~\citep{huang2021snapmix}.  The construction as presented is specific to
classification; extending the tile-scoring step to detection or segmentation,
where saliency is defined per region or per pixel, is left to future work.
